%==============================================================================
% فصل پنجم: بحث و نتیجه‌گیری ـ نسخه بازنویسی‌شده
% بازنویسی بر اساس فصل‌های دوم تا چهارم و کتابنامه ارسالی
%==============================================================================
\ChapterStart
\chapter{بحث و نتیجه‌گیری}
\label{ch:discussion}

\section{بحث}
\label{sec:discussion}
پژوهش حاضر با هدف بررسی ارزش افزوده خرده‌یادگیری در کنار آموزش کارگاهی خودمراقبتی در افراد مبتلا به پرفشاری خون انجام شد. از آنجا که هر دو گروه، برنامه کارگاهی حضوری مشابهی شامل ارائه کوتاه، پرسش و پاسخ، بحث گروهی و تمرین عملی دریافت کردند و تفاوت اصلی دو گروه در دریافت یا عدم دریافت میکروویدیوهای آموزشی بود، تفسیر یافته‌ها باید بر این پرسش متمرکز باشد که آیا تکرار و دسترسی پیوسته به محتوای کوتاه دیجیتال می‌تواند فراتر از اثر خود کارگاه، پیامدهای خودمراقبتی و بالینی را بهبود دهد. برآیند یافته‌ها نشان داد که پاسخ این پرسش برای همه پیامدها یکسان نیست: مزیت افزوده روش تلفیقی در عملکرد کلی خودمراقبتی مشاهده شد و این مزیت بیشتر در حیطه‌های اطلاعات، آگاهی و کنترل نمود یافت؛ در مقابل، در نگرش و رفتارها برتری مشخصی برای خرده‌یادگیری به دست نیامد. همچنین هر دو گروه از نظر فشار خون روند بهبود نشان دادند، اما شواهد کافی برای برتری روش تلفیقی نسبت به کارگاه در میزان تغییر فشار خون وجود نداشت.
\par
این الگو از نظر نظری معنادار است، زیرا پیامدهای آموزشی و شناختی به مداخله نزدیک‌ترند و می‌توانند در فاصله زمانی کوتاه‌تری به تکرار پیام، مرور محتوا و یادآوری پاسخ دهند؛ اما تغییر رفتار پایدار و سپس تغییر یک شاخص زیستی مانند فشار خون به حلقه‌های بیشتری از زنجیره اثر نیاز دارد. بیمار باید علاوه بر دانستن، بتواند آموخته‌ها را در شرایط واقعی زندگی اجرا کند، موانع را مدیریت کند، رفتار را تداوم دهد و در نهایت این تغییر رفتاری به اندازه‌ای پایدار باشد که در پیامد بالینی منعکس شود. این برداشت با جمع‌بندی مرور متون فصل دوم سازگار است که در آن اثر خرده‌یادگیری بر پیامدهای شناختی باثبات‌تر از اثر آن بر رفتار گزارش شده و برای تغییر رفتاری بر نقش تمرین، حمایت، بازخورد و پیگیری تأکید شده است \cite{wang2020microlearning}.

\section{تفسیر یافته مربوط به عملکرد کلی خودمراقبتی}
\label{sec:discussion-total-selfcare}
بهبود عملکرد خودمراقبتی در هر دو گروه نشان می‌دهد که خود آموزش کارگاهی، مستقل از جزء دیجیتال، ظرفیت قابل توجهی برای تقویت مراقبت از خود در افراد مبتلا به پرفشاری خون داشته است. این نتیجه با ماهیت برنامه اجراشده همخوان است؛ زیرا کارگاه صرفا انتقال اطلاعات نبود و شرکت‌کنندگان امکان پرسش، گفت‌وگو، مشاهده و تمرین مهارت‌هایی مانند اندازه‌گیری فشار خون را داشتند. در آموزش بزرگسالان، ارتباط مستقیم محتوا با مسئله واقعی زندگی، مشارکت فعال و بازخورد فوری از عناصر مهم یادگیری به شمار می‌روند و آموزش کارگاهی می‌تواند این عناصر را هم‌زمان فراهم کند \cite{knowles1980andragogy}.
\par
این بخش از یافته با نتایج مطالعات داخلی مرتبط با پرفشاری خون همسو است. مداخلات مبتنی بر الگوی \lr{PRECEDE-PROCEED} به بهبود رفتارهای خودمراقبتی منجر شده‌اند \cite{rakhshani2024precede,abedi2020precede} و آموزش مبتنی بر بحث گروهی نیز با بهبود رفتارهای خودمراقبتی همراه بوده است \cite{nasirian2019group}. وجه مشترک این مطالعات با پژوهش حاضر، عبور از آموزش یک‌سویه و استفاده از ساختار آموزشی هدفمند، مشارکت بیمار و پیوند دادن آموزش با رفتارهای روزمره است. بااین‌حال، تفاوت مهم آن است که بیشتر این پژوهش‌ها اثربخشی یک برنامه آموزشی را در برابر وضعیت معمول یا پیش از مداخله بررسی کرده‌اند، در حالی که پژوهش حاضر دو مداخله فعال را با یکدیگر مقایسه کرده است. بنابراین، یافته پژوهش حاضر علاوه بر تأیید سودمندی آموزش ساختاریافته، نشان می‌دهد که شیوه ارائه و تداوم دسترسی به محتوا نیز می‌تواند بر میزان بهبود خودمراقبتی اثر بگذارد.
\par
تریسون و همکاران (2008) بیان داشتند که برتری عملکرد کلی خودمراقبتی در گروه تلفیقی را می‌توان با نقش تقویتی خرده‌یادگیری تبیین کرد\cite{garrison2008blended}. کارگاه چهارچوب مفهومی و مهارتی اولیه را ایجاد می‌کرد، اما میکروویدیوها پیام‌های منتخب را در فاصله میان جلسات دوباره در دسترس بیمار قرار می‌دادند. کوتاه و متمرکز بودن هر واحد آموزشی، امکان مرور بدون نیاز به بازگشت به کل محتوای جلسه را فراهم می‌کرد و ارسال پیوسته محتوا می‌توانست به‌صورت یک محرک یادآور در محیط واقعی زندگی عمل کند. از این منظر، مزیت روش تلفیقی لزوما ناشی از «اطلاعات بیشتر» نیست، بلکه می‌تواند ناشی از «دسترسی مکرر و زمان‌بندی‌شده به همان پیام‌های اصلی» باشد. این تبیین با منطق آموزش ترکیبی، که نقش آموزش حضوری را در تعامل و تمرین و نقش بخش دیجیتال را در مرور و استمرار یادگیری مکمل یکدیگر می‌داند، سازگار است \cite{monib2024microlearning}.
\par
شواهد مطالعات آموزش ترکیبی و خرده‌یادگیری نیز از این تفسیر حمایت می‌کنند. در بیماران مبتلا به دیابت نوع دو، آموزش ترکیبی با بهبود دانش و رفتارهای خودمراقبتی همراه بوده است \cite{ghiyasvandian2023blended} و خرده‌یادگیری مبتنی بر رسانه اجتماعی در مقایسه با کارگاه متعارف، بهبود بیشتری در دانش، خودکارآمدی و رفتارهای خودمراقبتی نشان داده است \cite{rahbar2024socialmedia}. همچنین، آموزش مبتنی بر خرده‌یادگیری در سالمندان با ارتقای سواد سلامت و پایبندی دارویی همراه بوده است \cite{poorcheraghi2025microlearning}. هرچند جامعه و بیماری این مطالعات با پژوهش حاضر یکسان نیست، جهت کلی شواهد با این نتیجه همسو است که واحدهای کوتاه، قابل مرور و در دسترس می‌توانند جزء حضوری آموزش را تقویت کنند. از سوی دیگر، مرور نظام‌مند خرده‌یادگیری نشان داده است که اثر این رویکرد بر پیامدهای شناختی منظم‌تر از اثر آن بر رفتار است \cite{wang2020microlearning}؛ الگویی که در تحلیل حیطه‌های خودمراقبتی پژوهش حاضر نیز قابل مشاهده است.
\par
از منظر نظریه خودمراقبتی اورم، آموزش زمانی می‌تواند به مراقبت مؤثر از خود کمک کند که آگاهی، توان انجام رفتار و آمادگی فرد برای عمل تقویت شود \cite{orem1995nursing}. در پژوهش حاضر، کارگاه می‌توانسته از طریق تعامل و تمرین، توانایی اولیه را ایجاد کند و خرده‌یادگیری از طریق تکرار و یادآوری، دسترسی به دانش و استمرار توجه به رفتارهای مورد انتظار را تقویت کرده باشد. بااین‌حال، این مکانیسم در پژوهش حاضر به‌طور مستقیم اندازه‌گیری نشده است؛ بنابراین، باید آن را یک تبیین نظری سازگار با طراحی مداخله دانست، نه یک مسیر میانجی‌گری اثبات‌شده.

\section{تفسیر حیطه‌های خودمراقبتی}
\label{sec:discussion-domains}
الگوی حیطه‌های خودمراقبتی نشان داد که ارزش افزوده خرده‌یادگیری در همه ابعاد یکسان نبوده است. این ناهمگونی از نظر آموزشی اهمیت دارد، زیرا نشان می‌دهد افزودن ویدئوهای کوتاه به یک کارگاه مشترک، بیش از آنکه همه ابعاد خودمراقبتی را به یک اندازه تغییر دهد، بر برخی مؤلفه‌های نزدیک‌تر به یادگیری و مدیریت ادراک‌شده اثر گذاشته است.

\subsection{اطلاعات و آگاهی}
\label{subsec:discussion-information-awareness}
برتری گروه تلفیقی در حیطه‌های اطلاعات و آگاهی با ماهیت خود خرده‌یادگیری بیشترین سازگاری را دارد. میکروویدیوهای پژوهش بر پیام‌های مشخصی مانند شناخت پرفشاری خون، روش صحیح اندازه‌گیری فشار خون، تغذیه، فعالیت بدنی، مصرف دارو و مرور نکات خودمراقبتی متمرکز بودند. تکرار این پیام‌ها در واحدهای کوتاه می‌تواند بازیابی اطلاعات را آسان‌تر کند و احتمال فراموشی مطالب ارائه‌شده در جلسات حضوری را کاهش دهد. در چنین حالتی بیمار هنگام مواجهه با یک موقعیت واقعی، برای مثال زمان اندازه‌گیری فشار خون یا انتخاب غذا، به جای اتکا به حافظه یک جلسه طولانی، به یک پیام محدود و روشن دسترسی دارد.
\par
این یافته با مرور نظام‌مند خرده‌یادگیری همسو است که پیامدهای شناختی را از پایدارترین حوزه‌های پاسخ به این روش معرفی کرده است \cite{wang2020microlearning}. همچنین با مطالعاتی که آموزش ترکیبی را با افزایش دانش بیماران مبتلا به دیابت همراه دانسته‌اند \cite{ghiyasvandian2023blended} و مطالعه خرده‌یادگیری مبتنی بر رسانه اجتماعی که بهبود دانش را گزارش کرده است \cite{rahbar2024socialmedia} همخوانی دارد. مطالعه انجام‌شده در سالمندان نیز با نشان دادن بهبود سواد سلامت پس از خرده‌یادگیری، به‌طور غیرمستقیم از این مسیر شناختی حمایت می‌کند \cite{poorcheraghi2025microlearning}. بنابراین، منطقی است که نخستین اثر قابل مشاهده افزودن خرده‌یادگیری در پژوهش حاضر در حوزه‌هایی ظاهر شده باشد که به دریافت، حفظ و بازیابی اطلاعات وابستگی بیشتری دارند.

\subsection{کنترل}
\label{subsec:discussion-control}
برتری روش تلفیقی در حیطه کنترل را می‌توان با پیوند میان آموزش عملی و امکان مرور مجدد توضیح داد. همه شرکت‌کنندگان پیش از ثبت خانگی فشار خون، روش صحیح اندازه‌گیری را به‌صورت عملی آموختند و توانایی اجرای آن در حضور پژوهشگر بررسی شد؛ بنابراین، هر دو گروه پایه مهارتی مشترکی داشتند. بااین‌حال، در گروه تلفیقی یکی از میکروویدیوها به‌طور مشخص روش صحیح اندازه‌گیری فشار خون در منزل را مرور می‌کرد. دسترسی دوباره به یک نمایش کوتاه و تصویری می‌توانسته یادآوری مراحل کار را در محیط منزل تسهیل کند و احساس تسلط بیمار بر پایش بیماری را تقویت نماید.
\par
این تفسیر با چارچوب مفهومی پژوهش همخوان است که خرده‌یادگیری را تقویت‌کننده تکرار و دسترسی مستمر به پیام‌های آموزشی در نظر گرفته بود. همچنین با نظریه رفتار برنامه‌ریزی‌شده سازگار است که کنترل رفتاری ادراک‌شده را یکی از تعیین‌کننده‌های نزدیک رفتار می‌داند \cite{ajzen1991tpb}. بااین‌حال، باید توجه داشت که در مرحله پایه، حیطه کنترل بین دو گروه کاملاً متوازن نبود و در تحلیل نهایی مقدار پایه به‌عنوان عامل تعدیل‌کننده در نظر گرفته شد. ازاین‌رو، نتیجه این حیطه بهتر است به‌عنوان مؤید الگوی کلی پژوهش تفسیر شود، نه به‌عنوان شاهد مستقل و قطعی بر اثر خرده‌یادگیری.

\subsection{نگرش}
\label{subsec:discussion-attitude}
در حیطه نگرش، افزودن خرده‌یادگیری مزیت مشخصی نسبت به کارگاه ایجاد نکرد. یک تبیین محتمل آن است که بخش مشترک مداخله، یعنی کارگاه تعاملی، خود برای پرداختن به نگرش‌ها و باورهای مرتبط با بیماری ظرفیت بالایی داشته است. نگرش صرفا با تکرار اطلاعات تغییر نمی‌کند؛ بلکه برداشت فرد از خطر بیماری، منافع و موانع رفتار، تجربه شخصی، هنجارهای اجتماعی و ارزیابی او از توان انجام رفتار در شکل‌گیری آن نقش دارند. مدل باورهای سلامتی و نظریه رفتار برنامه‌ریزی‌شده نیز بر همین سازوکارها تأکید دارند \cite{rosenstock1974hbm,ajzen1991tpb}. بنابراین، گفت‌وگو، طرح نگرانی‌ها و تبادل تجربه در کارگاه ممکن است بخش مهمی از ظرفیت تغییر نگرش را در هر دو گروه فعال کرده باشد و تکرار ویدئویی، در بازه کوتاه پیگیری، اثر افزوده قابل تشخیصی ایجاد نکرده باشد.
\par
این نتیجه با مطالعه بحث گروهی در پرفشاری خون که بهبود مؤلفه‌های آموزشی و رفتاری را گزارش کرده است \cite{nasirian2019group} و نیز پژوهش آموزش گروهی در زنان که بهبود آگاهی و نگرش را نشان داده است \cite{ramazankhani2017group} قابل فهم است؛ زیرا در هر دو، تعامل و گفت‌وگو جزء اصلی آموزش بوده است. در پژوهش حاضر نیز هر دو گروه از چنین مؤلفه‌ای برخوردار بودند. بنابراین، نبود تفاوت بین دو گروه در نگرش را می‌توان نه به معنای بی‌اثر بودن آموزش، بلکه به معنای نبود شواهد کافی برای ارزش افزوده جزء ویدئویی بر اثر مشترک کارگاه تفسیر کرد.

\subsection{رفتارها}
\label{subsec:discussion-behaviors}
در حیطه رفتارها نیز مزیت افزوده معناداری برای خرده‌یادگیری مشاهده نشد. این یافته، در کنار بهبود بیشتر اطلاعات و آگاهی، فاصله میان «دانستن» و «انجام دادن» را برجسته می‌کند. رفتارهای خودمراقبتی پرفشاری خون، مانند مصرف منظم دارو، کاهش نمک، فعالیت بدنی و پایش مستمر فشار خون، رفتارهای تکرارشونده‌ای هستند که باید در روال روزمره تثبیت شوند. اجرای آنها تنها به یادآوری توصیه وابسته نیست و عوامل فردی و محیطی مانند انگیزه، خودکارآمدی، عادت‌های پیشین، حمایت خانواده، دسترسی به منابع و امکان دریافت بازخورد نیز بر تداوم رفتار اثر دارند. فصل دوم نیز بر چندعاملی بودن رفتار خودمراقبتی و نقش عوامل فردی، اجتماعی و سیستمی تأکید کرده است.
\par
این الگو با مرور نظام‌مند خرده‌یادگیری همخوان است؛ در آن مرور، بهبود پیامدهای شناختی نسبت به تغییر رفتاری باثبات‌تر گزارش شد و نتیجه‌گیری شد که محتوای کوتاه به تنهایی برای ایجاد تغییر رفتاری پایدار کافی نیست \cite{wang2020microlearning}. در مقابل، مطالعه حمایت از خودمدیریتی در بیماران مبتلا به پرفشاری خون، که علاوه بر آموزش بر حمایت از تصمیم‌گیری و مدیریت خودمراقبتی تمرکز داشت، بهبود پایبندی و مدیریت خودمراقبتی را گزارش کرده است \cite{kurt2022selfmgmt}. این تفاوت می‌تواند نشان دهد که برای عبور از یادگیری به رفتار، افزودن مؤلفه‌هایی مانند هدف‌گذاری فردی، خودپایشی، حل مسئله و بازخورد ممکن است اهمیت بیشتری از صرف تکرار محتوای آموزشی داشته باشد. این توضیح در پژوهش حاضر مستقیماً آزمون نشده است و باید به‌عنوان فرضیه‌ای برای تبیین یافته و طراحی مطالعات بعدی در نظر گرفته شود.

\section{تفسیر یافته‌های فشار خون سیستولیک و دیاستولیک}
\label{sec:discussion-bp}
کاهش فشار خون در هر دو گروه، از نظر جهت، با انتظار نظری از یک برنامه آموزش خودمراقبتی سازگار است. محتوای مشترک کارگاه رفتارهایی را هدف قرار می‌داد که با مدیریت پرفشاری خون ارتباط مستقیم دارند؛ از جمله مصرف دارو مطابق تجویز، کاهش نمک و چربی، فعالیت بدنی و پایش خانگی فشار خون. بنابراین، بهبود فشار خون در هر دو گروه می‌تواند با فعال شدن مجموعه‌ای از رفتارهای مراقبتی پس از آموزش سازگار باشد. این جهت یافته با مطالعاتی که پس از مداخلات آموزشی یا حمایتی در بیماران مبتلا به پرفشاری خون کاهش فشار خون را گزارش کرده‌اند، همسو است؛ از جمله مداخله مبتنی بر \lr{PRECEDE-PROCEED} \cite{rakhshani2024precede}، آموزش به روش بحث گروهی \cite{nasirian2019group} و برنامه حمایت از خودمدیریتی \cite{kurt2022selfmgmt}.
\par
بااین‌حال، نکته اصلی پژوهش حاضر این نیست که آیا آموزش به‌طور کلی با کاهش فشار خون همراه بوده است، بلکه این است که آیا افزودن خرده‌یادگیری نسبت به یک کارگاه فعال، کاهش بیشتری ایجاد کرده است. بر اساس مقایسه تغییرات بین دو گروه، چنین برتری‌ای تأیید نشد. این نتیجه الزاماً با مطالعات یادشده متناقض نیست، زیرا سؤال مقایسه‌ای آنها متفاوت بوده است. در بسیاری از مطالعات پیشین، یک مداخله آموزشی با مراقبت معمول یا وضعیت قبل از مداخله مقایسه شده است؛ در حالی که در پژوهش حاضر، گروه مقایسه نیز چهار جلسه آموزش کارگاهی فعال دریافت کرده بود. در نتیجه، هر دو گروه در معرض مؤلفه‌هایی قرار داشتند که می‌توانند بر رفتارهای مؤثر بر فشار خون اثر بگذارند و فضای کمتری برای آشکار شدن اثر افزوده میکروویدیوها باقی می‌ماند.
\par
تفاوت الگوی فشار خون با خودمراقبتی را می‌توان از منظر فاصله پیامد از مداخله نیز توضیح داد. اطلاعات و آگاهی مستقیماً در معرض آموزش قرار می‌گیرند، اما فشار خون یک پیامد زیستی چندعاملی است. برای آنکه تفاوت آموزشی به تفاوت پایدار فشار خون تبدیل شود، تغییرات رفتاری باید به اندازه کافی اجرا و تثبیت شوند و هم‌زمان عوامل دیگری مانند درمان دارویی، دریافت نمک، فعالیت بدنی، وزن، خواب و تنش روانی نیز بر نتیجه اثر می‌گذارند. پژوهش حاضر این عوامل را به‌عنوان میانجی یا مخدوش‌کننده اندازه‌گیری و تحلیل نکرده است؛ ازاین‌رو، نمی‌توان مسیر دقیق تغییر فشار خون را به یک سازوکار مشخص نسبت داد. افزون بر این، پیگیری یک‌ماهه برای مشاهده اثرات کوتاه‌مدت مناسب است، اما ممکن است برای ایجاد فاصله روشن بین دو برنامه فعال در یک پیامد فیزیولوژیک کافی نباشد.
\par
نحوه سنجش فشار خون در پژوهش حاضر یک نکته مثبت در تفسیر این یافته است. فشار خون پایه و پیگیری بر مبنای ثبت خانگی مکرر طی یک هفته تعیین شد و شرکت‌کنندگان پیش از ثبت، آموزش عملی دیدند. این روش نسبت به اتکا به یک اندازه‌گیری منفرد، امکان نمایش بهتر وضعیت معمول بیمار را فراهم می‌کند. با وجود این، چون پژوهش گروه مراقبت معمول نداشت و اطلاعات ارائه‌شده در فصل یافته‌ها امکان تفکیک اثر آموزش از تمام تغییرات هم‌زمان درمانی و زمانی را فراهم نمی‌کند، کاهش مطلق فشار خون در هر دو گروه نباید به‌طور کامل به مداخله آموزشی نسبت داده شود. همچنین، معنادار نبودن تفاوت بین دو گروه به معنای اثبات برابری یا هم‌ارزی دو روش نیست؛ بلکه تنها نشان می‌دهد در این مطالعه شواهد کافی برای برتری یکی از دو روش در میزان تغییر فشار خون به دست نیامده است.

\section{تبیین یکپارچه یافته‌ها}
\label{sec:integrated-interpretation}
در کنار هم قرار دادن یافته‌های خودمراقبتی و فشار خون، یک الگوی مرحله‌ای را پیشنهاد می‌کند. نخست، آموزش کارگاهی مشترک در هر دو گروه می‌تواند فهم اولیه، فرصت تمرین و آمادگی برای خودمراقبتی را ایجاد کرده باشد. دوم، خرده‌یادگیری با تکرار پیام‌ها و امکان مرور در فاصله میان جلسات، برخی پیامدهای نزدیک‌تر به یادگیری مانند اطلاعات و آگاهی و تا حدی احساس کنترل را بیشتر تقویت کرده است. سوم، این مزیت شناختی در بازه کوتاه مطالعه به برتری مشخص در همه رفتارهای خودمراقبتی تبدیل نشده است. چهارم، در سطح پیامد زیستی نیز تفاوت افزوده‌ای برای فشار خون آشکار نشده است. چنین الگویی با چارچوب مفهومی فصل دوم سازگار است و نشان می‌دهد مسیر اثر آموزش از یادگیری تا پیامد بالینی خطی و فوری نیست.
\par
از منظر مدل خودمراقبتی در بیماری مزمن، مدیریت موفق بیماری مستلزم تداوم رفتارهای نگهدارنده، پاسخ مناسب به تغییر وضعیت و اطمینان فرد به توان مدیریت بیماری است \cite{riegel2012middlerange}. خرده‌یادگیری می‌تواند بخشی از این فرایند را از طریق یادآوری و دسترس‌پذیری پشتیبانی کند، اما برای ایجاد تغییر رفتاری پایدار احتمالاً باید با سازوکارهای دیگری مانند بازخورد، خودپایشی و حل مسئله همراه شود. بنابراین، یافته‌های پژوهش حاضر بیشتر از نقش «مکمل» خرده‌یادگیری حمایت می‌کنند تا نقش «جایگزین» آن برای آموزش تعاملی حضوری.

\section{پیامدهای آموزشی و کاربردی}
\label{sec:implications}
وانگ و همکاران (2020) نشان دادند که خرده‌یادگیری می‌تواند به‌عنوان لایه تقویتی پس از و یا میان جلسات کارگاهی مورد استفاده قرار گیرد، به‌ویژه برای مطالبی که نیازمند مرور مکرر هستند\cite{wang2020microlearning}. موضوعاتی مانند روش اندازه‌گیری فشار خون، نکات تغذیه‌ای، مصرف دارو و یادآوری فعالیت بدنی قابلیت تبدیل شدن به واحدهای کوتاه و مستقل را دارند. مزیت اصلی چنین واحدهایی در دسترسی سریع و امکان مرور مجدد است، نه در جایگزین کردن ارتباط درمانی و آموزش حضوری.
\par
در طراحی برنامه‌های آینده بهتر است هر واحد خرده‌یادگیری یک هدف روشن و قابل اجرا داشته باشد و مستقیما به تکلیفی در زندگی روزمره بیمار متصل شود. برای مثال، پس از یک ویدئوی مربوط به پایش فشار خون، ثبت فشار خون یا مرور خطاهای اندازه‌گیری می‌تواند به‌عنوان تکلیف پیگیری شود. برای حوزه‌هایی که در پژوهش حاضر برتری افزوده مشخصی نشان ندادند، به‌ویژه رفتار و نگرش، افزودن مؤلفه‌های تغییر رفتار مانند هدف‌گذاری فردی، ثبت عملکرد، بازخورد، حل موانع و درگیر کردن حمایت خانواده می‌تواند منطقی‌تر از افزایش صرف حجم محتوای ویدئویی باشد. این پیشنهاد با چارچوب‌های رفتاری مطرح‌شده در فصل دوم و با این مشاهده سازگار است که دانش به تنهایی برای تغییر رفتار کافی نیست \cite{rosenstock1974hbm,ajzen1991tpb}.
\par
در محیط بالینی نیز ارزیابی موفقیت برنامه آموزشی بهتر است در چند سطح انجام شود: یادگیری و آگاهی، اجرای واقعی رفتارهای خودمراقبتی و پیامد بالینی. یافته‌های این مطالعه نشان می‌دهد بهبود در سطح نخست و حتی نمره کلی خودمراقبتی لزوما به معنای ایجاد مزیت فوری در فشار خون نیست. بنابراین، پایش هم‌زمان رفتارهای عینی و شاخص‌های بالینی، به‌ویژه در پیگیری‌های طولانی‌تر، برای داوری درباره اثربخشی کامل برنامه ضروری است.

\section{نقاط قوت پژوهش}
\label{sec:strengths}
یکی از نقاط قوت پژوهش، استفاده از طرح کارآزمایی بالینی تصادفی‌شده با گروه مقایسه فعال بود. یکسان بودن آموزش کارگاهی در دو گروه و تفاوت داشتن آنها در جزء خرده‌یادگیری، امکان تفسیر روشن‌تری از ارزش افزوده محتوای کوتاه دیجیتال فراهم کرد. این ویژگی از نظر پاسخ به شکاف مطرح‌شده در فصل دوم اهمیت دارد، زیرا بسیاری از مطالعات پیشین یک برنامه آموزشی را با مراقبت معمول مقایسه کرده بودند و سهم افزوده خرده‌یادگیری را جداگانه نمی‌سنجیدند.
\par
نقطه قوت دیگر، طراحی ساختاریافته آموزش بر اساس الگوی \lr{ADDIE} و هماهنگی محتوای میکروویدیوها با محورهای اصلی خودمراقبتی بود \cite{branch2009addie}. همچنین، برای سنجش فشار خون از ثبت خانگی مکرر و آموزش عملی روش اندازه‌گیری استفاده شد و عملکرد خودمراقبتی با نسخه فارسی ابزار اختصاصی پرفشاری خون سنجیده شد که روایی و پایایی آن پیش‌تر بررسی شده بود \cite{ghaneigheshlagh2019persian}. این ویژگی‌ها به هم‌راستایی مداخله، ابزار و پیامدهای مورد ارزیابی کمک کرده‌اند.

\section{محدودیت‌های پژوهش}
\label{sec:limitations}
نخست، پیگیری پژوهش به یک ماه پس از پایان مداخله محدود بود. این فاصله برای بررسی پاسخ کوتاه‌مدت مناسب است، اما امکان قضاوت درباره ماندگاری یادگیری، تثبیت رفتارهای خودمراقبتی و پایداری تفاوت‌های فشار خون را فراهم نمی‌کند. این محدودیت به‌ویژه برای تفسیر رفتار و فشار خون مهم است، زیرا این پیامدها ممکن است برای تثبیت به زمان بیشتری نیاز داشته باشند.
\par
دوم، از ۱۱۰ شرکت‌کننده واردشده به مطالعه، داده پس‌آزمون ۱۶ نفر در دسترس نبود و تحلیل گزارش‌شده بر داده‌های کامل انجام شد. فایل‌های در دسترس اطلاعاتی درباره علت این فقدان داده ارائه نمی‌کنند؛ بنابراین، نمی‌توان مشخص کرد که ریزش تا چه اندازه با ویژگی شرکت‌کنندگان یا پاسخ به مداخله ارتباط داشته است. این موضوع می‌تواند بر تعمیم و برآورد دقیق اثر مداخله محدودیت ایجاد کند.
\par
سوم، پژوهش فاقد گروه مراقبت معمول و نیز فاقد گروه خرده‌یادگیریِ بدون کارگاه بود. در نتیجه، طراحی حاضر برای برآورد ارزش افزوده خرده‌یادگیری نسبت به کارگاه مناسب است، اما سهم مطلق کارگاه در مقایسه با روند طبیعی مراقبت و نیز اثر مستقل خرده‌یادگیری را مشخص نمی‌کند. این نکته در تفسیر کاهش فشار خون اهمیت ویژه دارد، زیرا هر دو گروه آموزش فعال دریافت کرده‌اند.
\par
چهارم، عملکرد خودمراقبتی با پرسشنامه خودگزارش‌شده سنجیده شد. هرچند ابزار مورد استفاده برای پرفشاری خون طراحی و نسخه فارسی آن اعتبارسنجی شده است \cite{han2014hscp,ghaneigheshlagh2019persian}، پاسخ‌های خودگزارش‌شده می‌توانند با یادآوری فرد و تمایل به ارائه پاسخ مطلوب تحت تأثیر قرار گیرند. استفاده هم‌زمان از شاخص‌های عینی رفتار، مانند ثبت واقعی مصرف دارو، فعالیت بدنی یا الگوی غذایی، می‌توانست تصویر کامل‌تری از تغییر رفتاری ارائه دهد.
\par
پنجم، معیارهای ورود شامل توانایی خواندن و نوشتن، دسترسی به تلفن هوشمند و اینترنت، توان استفاده از محتوای ارسالی و دامنه سنی مشخص بود و نمونه از یک مرکز درمانی جذب شد. بنابراین، تعمیم یافته‌ها به بیمارانی با سواد سلامت یا سواد دیجیتال محدود، افراد فاقد دسترسی پایدار به اینترنت، گروه‌های سنی خارج از دامنه مطالعه و محیط‌های درمانی دیگر باید با احتیاط انجام شود.
\par
ششم، در فصل روش‌ها پایش مشاهده میکروویدیوها و دریافت تأیید مشاهده به‌عنوان بخشی از ارزشیابی فرایندی ذکر شده است، اما در فصل یافته‌های در دسترس شاخص کمیِ میزان مواجهه، تعداد دفعات مرور یا الگوی مشاهده ویدئوها گزارش نشده است. ازاین‌رو، نمی‌توان بررسی کرد که آیا میزان استفاده واقعی از خرده‌یادگیری با شدت پاسخ خودمراقبتی ارتباط داشته است یا خیر.
\par
هفتم، برای حیطه‌های متعدد خودمراقبتی چند مقایسه آماری انجام شده است و در بخش روش‌های آماریِ فایل‌های ارسالی، روش مشخصی برای اصلاح تعدد آزمون‌ها گزارش نشده است. بنابراین، نتایج حیطه‌ای بهتر است به‌عنوان تحلیل‌های تکمیلی و تبیین‌کننده الگوی پیامد اصلی در نظر گرفته شوند و برای تأیید مستقل هر حیطه، تکرار یافته‌ها در نمونه‌های دیگر ضروری است.

\section{پیشنهادها برای پژوهش‌های آینده}
\label{sec:future-research}
پیشنهاد می‌شود مطالعات آینده با پیگیری‌های طولانی‌تر انجام شوند تا مشخص شود مزیت کوتاه‌مدت خرده‌یادگیری در خودمراقبتی تا چه اندازه پایدار می‌ماند و آیا در بازه زمانی طولانی‌تر به تفاوت رفتاری و بالینی روشن‌تری منجر می‌شود. استفاده از سنجش‌های چندمرحله‌ای می‌تواند ترتیب زمانی تغییرات را نیز آشکار کند؛ برای مثال ابتدا اطلاعات و آگاهی، سپس رفتار خودمراقبتی و در نهایت فشار خون بررسی شود.
\par
بهتر است در طراحی‌های بعدی چهار بازوی مراقبت معمول، کارگاه، خرده‌یادگیری و روش تلفیقی در نظر گرفته شود تا اثر مستقل هر جزء و اثر ترکیب آنها از یکدیگر تفکیک شود. در صورت تمرکز بر پیامد فشار خون، ثبت تغییرات دارویی و شاخص‌های رفتاری مرتبط مانند مصرف نمک، فعالیت بدنی، وزن و پایبندی به درمان می‌تواند به تبیین مسیر اثر کمک کند.
\par
برای بررسی رفتار واقعی، استفاده از پرسشنامه در کنار شاخص‌های عینی توصیه می‌شود. ثبت الکترونیکی یا ساختاریافته مصرف دارو، دفعات پایش فشار خون، فعالیت بدنی و برخی شاخص‌های تغذیه‌ای می‌تواند فاصله میان ادراک خودمراقبتی و رفتار واقعی را بهتر نشان دهد. همچنین، ثبت خودکار تعداد مشاهده، مدت مشاهده و مرور مجدد هر میکروویدیو امکان تحلیل رابطه میان میزان مواجهه با آموزش و پاسخ به مداخله را فراهم می‌کند.
\par
با توجه به نبود مزیت افزوده روشن در نگرش و رفتارها، مطالعات آینده می‌توانند خرده‌یادگیری را با راهبردهای فعال تغییر رفتار ترکیب کنند؛ از جمله هدف‌گذاری فردی، خودپایشی، بازخورد شخصی، حل مسئله و حمایت خانواده یا همتایان. بررسی میانجی‌ها نیز می‌تواند روشن کند که آیا افزایش اطلاعات و آگاهی از طریق تقویت کنترل ادراک‌شده، خودکارآمدی یا رفتارهای مشخص به پیامد بالینی منجر می‌شود \cite{riegel2012middlerange,ajzen1991tpb}.
\par
در نهایت، لازم است امکان‌پذیری اجرای روش تلفیقی در گروه‌های دارای دسترسی یا سواد دیجیتال کمتر نیز بررسی شود. ارزیابی پذیرش بیمار، سهولت استفاده، هزینه تولید و ارسال محتوا و گزینه‌های جایگزین برای افرادی که تلفن هوشمند یا اینترنت پایدار ندارند، برای تصمیم‌گیری درباره اجرای گسترده این روش اهمیت دارد.

\section{نتیجه‌گیری}
\label{sec:conclusion}
یافته‌های پژوهش نشان می‌دهد آموزش کارگاهی برای افراد مبتلا به پرفشاری خون، در قالبی تعاملی و مهارت‌محور، با بهبود عملکرد خودمراقبتی همراه است و افزودن خرده‌یادگیری می‌تواند این اثر را در کوتاه‌مدت تقویت کند. مزیت افزوده روش تلفیقی بیشتر در حیطه‌هایی ظاهر شد که به دریافت، یادآوری و مدیریت اطلاعات نزدیک‌ترند؛ یعنی اطلاعات، آگاهی و کنترل. در مقابل، برای نگرش و رفتارها شواهدی از برتری افزوده خرده‌یادگیری نسبت به کارگاه به دست نیامد. این الگو با شواهد موجود درباره خرده‌یادگیری سازگار است و نشان می‌دهد محتوای کوتاه و تکرارشونده در تثبیت یادگیری ظرفیت بیشتری دارد، اما تبدیل یادگیری به رفتار پایدار احتمالاً به راهبردهای تکمیلی نیاز دارد.
\par
از نظر فشار خون نیز هر دو گروه روند بهبود داشتند، اما روش تلفیقی نسبت به کارگاه به‌تنهایی مزیت افزوده مشخصی در میزان تغییر فشار خون نشان نداد. بنابراین، یافته‌های مطالعه از به‌کارگیری خرده‌یادگیری به‌عنوان مکمل آموزش کارگاهی برای تقویت خودمراقبتی حمایت می‌کند، اما برای ادعای برتری آن در پیامد بالینی فشار خون کافی نیست. بر این اساس، کاربرد مناسب خرده‌یادگیری در این زمینه، جایگزین کردن آموزش حضوری نیست؛ بلکه ایجاد یک لایه کوتاه، قابل مرور و تکرارشونده برای حفظ ارتباط آموزشی میان جلسات و تقویت پیام‌های کلیدی است. برای تعیین پایداری اثر و تبدیل آن به رفتار و پیامد بالینی، پیگیری طولانی‌تر و مداخلات رفتاری تکمیلی مورد نیاز است.

